\section{The Second Document}
\label{sec:ch07-the-second-document}

\index{composition!a graph of one hides everything|(}%
A graph of one subgraph composes, and that is the problem with it. Every
question composition exists to answer is a question about two services
disagreeing, so a graph with nothing to disagree with passes by having nothing
to check. The config in
section~\ref{sec:ch07-what-the-composer-wrote} is real and a router will serve
it, and it tells you nothing about whether this subgraph can live beside
another one.

\index{composition!reads documents, not services}%
Which you can find out now, because of the first fact in this chapter. The
composer reads files. It does not care whether a service exists behind one, so
a second document is enough to ask the question, and a second document is
cheaper than a second service by about a chapter.

\index{speakers.graphql@\code{speakers.graphql}!is written by hand}%
Here is the one this book uses. It is not exported by anything. Nobody is
running a Speakers service and chapter~\ref{ch:second-third-subgraph} is where
one gets built; this is a schema I wrote by hand to find out what will happen
when it does:

\begin{graphqlsdl}
schema
  @link(
    url: "https://specs.apollo.dev/federation/v2.6"
    import: ["@key", "@shareable", "FieldSet"]
  ) {
  query: Query
}

type Query {
  node(id: ID!): Node @shareable
  nodes(ids: [ID!]!): [Node]! @shareable
}

interface Node {
  id: ID!
}

type Speaker implements Node @key(fields: "id") {
  id: ID!
  name: String!
  bio: String
}

scalar FieldSet
\end{graphqlsdl}

\index{speakers.graphql@\code{speakers.graphql}!why it declares node}%
Twenty-four lines, and two of them are doing all the work.
Chapter~\ref{ch:schema-design} switched global object identification on, and
every Hot Chocolate service this book builds will have it on for the same
reason this one does. So every subgraph in this graph exports \code{node} and
\code{nodes}, and this document says so. Everything else in it is the Speakers
subgraph as chapter~\ref{ch:federation-model} drew it.

\index{graph.yaml@\code{graph.yaml}!a second graph file}%
A second input file names both:

\begin{minted}{yaml}
version: 1
subgraphs:
  - name: sessions
    routing_url: http://localhost:5001/graphql
    schema:
      file: ../src/Sessions/schema.graphql
  - name: speakers
    routing_url: http://localhost:5002/graphql
    schema:
      file: ./speakers.graphql
\end{minted}

\subsection{The Bill, Presented}

\index{composition!the collision arrives}%
Compose that and the answer is the one chapter~\ref{ch:federation-model} said
was coming:

\begin{minted}[fontsize=\footnotesize,breakanywhere]{text}
The Object "Query" defines the same fields in multiple subgraphs without the "@shareable" directive:
 The field "node" is defined and declared "@shareable" in the following subgraph: "speakers".
 However, it is not declared "@shareable" in the following subgraph: "sessions".
 The field "nodes" is defined and declared "@shareable" in the following subgraph: "speakers".
 However, it is not declared "@shareable" in the following subgraph: "sessions".
The Object "Speaker" defines the same fields in multiple subgraphs without the "@shareable" directive:
 The field "name" is defined in the following subgraphs: "sessions", "speakers".
 However, it is not declared "@shareable" in any of them.
 The field "bio" is defined in the following subgraphs: "sessions", "speakers".
 However, it is not declared "@shareable" in any of them.
\end{minted}

\index{shareable@\code{"@shareable}!two forms of one message}%
Read the two halves against each other, because the message takes a different
shape in each and the difference is useful. For \code{Query} it names the
subgraph that is missing the directive, because the other one has it. For
\code{Speaker} it says neither of them has it. Same rule, two reports, and the
first is the one you want: it tells you where to go.

\index{ownership!two collisions, opposite fixes}%
What the composer cannot tell you is that these two collisions want opposite
answers. \code{node} is a field two services genuinely both resolve, and every
service in this graph will resolve it, so the fix is to say so in both.
\code{Speaker.name} is a field two services both claim and only one should own.
Sessions still stores the speakers, which chapter~\ref{ch:federation-model}
said would stay true until there was somewhere to move the rows to, and moving
them is chapter~\ref{ch:second-third-subgraph}'s work.

\index{composition!says what collides, not who should yield}%
That distinction is not in the error and could not be. Composition can see that
two documents claim one field. Which of the two ought to give it up is a
question about who owns the data, and chapter~\ref{ch:ownership} is a whole
chapter on how badly that question goes when nobody answers it deliberately.

\subsection{Where the Declaration Goes}

\index{shareable@\code{"@shareable}!there is no method to put it on}%
So \code{Query.node} needs \code{@shareable} on this side. The obvious move is
an attribute, and the obvious place is the method the field comes from. There
is no method. \code{node} and \code{nodes} are generated by
\code{AddGlobalObjectIdentification()}, and no class in this service declares
either of them. That rules out the idiom the rest of this book uses before the
question is properly asked.

\index{Shareable@\code{[Shareable]}!on the Query class}%
The attribute does have somewhere to go, one level up. \code{[Shareable]} is
legal on a class, and \code{Query} is a class, so putting it there compiles and
exports this:

\begin{graphqlsdl}
type Query @shareable {
\end{graphqlsdl}

\index{shareable@\code{"@shareable}!on an object means every field}%
That clears the collision. It also clears things that were not collisions.
\code{@shareable} on an object means every field of that object, so
\code{sessions} and \code{sessionById} are declared shareable now as well, and
what that permits is worth measuring rather than imagining. So I measured it:
against a subgraph whose \code{Query} is shareable as a type, a second document
that declares its own \code{sessionById} composes. Not with a softer warning.
Exit zero, a router execution config written, nothing said. A field only this
service can answer is now a field any subgraph may claim, and the composer has
been told not to mind.

\index{ownership!a check switched off is not a check that warns}%
In a book with a chapter on the subgraph that broke everyone and another on who
owns what, that is the wrong thing to switch off to save some lines.

\index{ObjectTypeExtension@\code{ObjectTypeExtension}!does not merge here}%
\code{ObjectTypeExtension} is the schema-building route, and it cannot reach
these fields at all. One naming \code{Query} and configuring
\code{Field("node")} does not merge with the generated field; it declares a
second field with the same name. Without an explicit type the schema fails to
build with an error about not being able to infer the type of
\code{Query.node}, and with the type supplied it fails with an error about
\code{node} and \code{nodes} being declared multiple times on \code{Query}. I
am describing those rather than printing them, because a service in either
state does not start and neither is a state this book ships.

\index{TypeInterceptor@\code{TypeInterceptor}!reaches a generated field}%
What does reach exactly the two fields is a type interceptor. It runs while the
schema is being built, sees the configuration of each type before that type is
finished, and can change it. This one is the whole file:

\begin{minted}{csharp}
using HotChocolate.Configuration;
using HotChocolate.Language;
using HotChocolate.Types.Descriptors.Configurations;

namespace Sessions;

/// <summary>
/// Applies @shareable to the two fields global object identification
/// generates, and to nothing else.
/// </summary>
internal sealed class ShareNodeFields : TypeInterceptor
{
    public override void OnBeforeCompleteType(
        ITypeCompletionContext context,
        TypeSystemConfiguration configuration)
    {
        if (configuration is not ObjectTypeConfiguration type
            || type.Name != "Query")
        {
            return;
        }

        foreach (var field in type.Fields)
        {
            if (field.Name is "node" or "nodes")
            {
                field.Directives.Add(
                    new DirectiveConfiguration(
                        new DirectiveNode("shareable")));
            }
        }
    }
}
\end{minted}

\index{TryAddTypeInterceptor@\code{TryAddTypeInterceptor}!registers it}%
Registering it is one more line on the builder, below the cost option from
section~\ref{sec:ch07-whose-directive-is-it}. Here is \code{Program.cs} in its
final state for this chapter, with both of the chapter's lines marked:

\begin{minted}[highlightlines={35,36}]{csharp}
using ChilliCream.Nitro.App;
using HotChocolate.AspNetCore;
using Microsoft.EntityFrameworkCore;
using Sessions;

var builder = WebApplication.CreateBuilder(args);

// EF Core reports every command through the logging pipeline as well,
// at Information and with a duration attached. StatementLog below is
// this service's instrument, so silence the other one rather than have
// every statement printed twice.
builder.Logging.AddFilter(
    "Microsoft.EntityFrameworkCore.Database.Command",
    LogLevel.Warning);

builder.Services.AddDbContextFactory<ConferenceContext>(options =>
{
    options.UseSqlite("Data Source=conference.db");

    if (builder.Environment.IsDevelopment())
    {
        options.AddInterceptors(new StatementLog());
    }
});

builder
    .AddGraphQL()
    .AddSessionsTypes()
    .RegisterDbContextFactory<ConferenceContext>()
    .AddQueryContext()
    .AddPagingArguments()
    .AddGlobalObjectIdentification()
    .AddMutationConventions()
    .AddApolloFederation()
    .ModifyCostOptions(options => options.ApplyCostDefaults = false)
    .TryAddTypeInterceptor<ShareNodeFields>();

var app = builder.Build();

using (var scope = app.Services.CreateScope())
{
    var factory = scope.ServiceProvider
        .GetRequiredService<IDbContextFactory<ConferenceContext>>();
    using var db = factory.CreateDbContext();
    db.Database.EnsureCreated();
}

app.MapGraphQL()
    .WithOptions(options =>
    {
        // Serve the Nitro build that shipped in the package. The
        // default, ServeMode.Latest, downloads the current one
        // from ChilliCream's CDN instead.
        options.Tool.ServeMode = ServeMode.Embedded;
    });

app.RunWithGraphQLCommands(args);
\end{minted}

\index{shareable@\code{"@shareable}!on two fields and nothing else}%
And the export says what it should:

\begin{graphqlsdl}
type Query {
  node("ID of the object." id: ID!): Node @shareable
  nodes("The list of node IDs." ids: [ID!]!): [Node]! @shareable
  ...
}
\end{graphqlsdl}

\index{TypeInterceptor@\code{TypeInterceptor}!is the right size for this}%
Thirty-three lines against one attribute is a bad trade on every measure except
the one that decides it, which is what the schema then says. The attribute
makes a promise about four fields in order to buy a promise about two. This
makes the promise about two. I would spend the lines again, and I would rather
Hot Chocolate made the question unnecessary by declaring these two the way it
already declares \code{PageInfo}.

\index{composition!what remains after the fix}%
Compose the two documents now and \code{Query} is gone from the error
entirely. What is left is \code{Speaker}, alone, which is the seam.

\begin{onhc14}
Everything in this chapter happens on~14 as well, and the failure comes first.
A 14.3.1 subgraph exports the same \code{@cost(weight: "10")} and the same
\code{@listSize} carrying a \code{slicingArgumentDefaultValue}, so it does not
compose either, and for a reason that is no more federation's there than here.

The first fix ports unchanged. \code{ModifyCostOptions} is on 14.3.1 under that
name, \code{CostOptions} carries the same nine properties, and
\code{ApplyCostDefaults} set false removes both directives exactly as it does
on~16.

The second ports after four renames, all of them names~16 changed rather than
mechanisms it added. The class body is otherwise the same file:

\begin{center}
\begin{tabular}{ll}
\toprule
On 16.6.1 & On 14.3.1 \\
\midrule
\code{TypeSystemConfiguration} & \code{DefinitionBase} \\
\code{ObjectTypeConfiguration} & \code{ObjectTypeDefinition} \\
\code{DirectiveConfiguration} & \code{DirectiveDefinition} \\
\code{Descriptors.Configurations} & \code{Descriptors.Definitions} \\
\bottomrule
\end{tabular}
\end{center}

So the override reads:

\begin{minted}[fontsize=\footnotesize]{csharp}
public override void OnBeforeCompleteType(
    ITypeCompletionContext context,
    DefinitionBase definition)
\end{minted}

and everything below it, including the \code{DirectiveNode("shareable")} the
directive is built from, is character for character what~16 uses.
Appendix~\ref{app:hc14} carries the class.

One difference is not this chapter's. The key table in the router execution
config has one row on this branch where~16 has two, because
chapter~\ref{ch:first-subgraph} established that \code{Session} never becomes
an entity here. That has a consequence worth knowing: run this chapter's
unreachable-field experiment against a 14 subgraph and you never reach the
resolvability pass at all. \code{Session} is not an entity, so both documents
merely declare \code{Session.id}, and the merge fails on shareability before
anything walks anywhere.
\end{onhc14}
\index{composition!a graph of one hides everything|)}%
